\section{Experiments}
\label{sec:exp}

We evaluate CFA with three questions in mind: how it compares with existing
adversarial training and robust distillation methods, whether its performance
generalizes across datasets and model architectures, and how its individual
design components contribute to the resulting clean--robust trade-off.


\subsection{Experimental Settings}

We evaluate CFA on CIFAR-10, CIFAR-100
\citep{krizhevsky2009learning}, and Tiny-ImageNet-200
\citep{le2015tiny}, using ResNet-18 as the default architecture.
We additionally evaluate WideResNet-34-10 to examine generalization across
model capacity.

Unless stated otherwise, the clean self-teacher is trained on natural images
for 200 epochs and then frozen. The student is initialized from the teacher,
with the inherited classifier kept fixed during adversarial training.
Following \Cref{sec:method}, the student is trained using clean feature
anchoring with sensitivity-matched perturbation radii (\(p=1\)).
The main CIFAR experiments use 150 training epochs and
\(\epsilon_{\mathrm{train}}=8/255\).
We use AdamW with a one-cycle learning-rate schedule and 10-step PGD,
together with weight averaging (WA) and adversarial weight perturbation
(AWP). Teacher statistics and dataset- and architecture-specific training
settings are reported in \Cref{tab:teachers,app:repro}.

We compare CFA with standard adversarial training and robust distillation
baselines using their respective training recipes. Distillation-based
baselines receive the same clean self-teacher where applicable.
We report clean accuracy and AutoAttack (AA) accuracy on the full test set at
\(\epsilon=8/255\), together with their arithmetic mean (Avg) and harmonic
mean (NRR). Additional evaluation details, attack diagnostics, and results
with alternative teachers are provided in
\Cref{app:repro,app:gradmask,tab:published}.


\subsection{Main Results}

\begin{table}[!htbp]
  \centering
  \footnotesize
  \setlength{\tabcolsep}{3.2pt}
  \caption{
  ResNet-18 results under \(\ell_\infty\) perturbations with
  \(\epsilon=8/255\). All values are measured in our framework.
  \emph{T.} denotes the teacher assumed by each method: none (--),
  a natural model (nat.), or the method's own weight average (self).
  The CIFAR-100 CFA row reports the mean of three student seeds;
  all other rows are individual runs.
  }
  \label{tab:main}
  \begin{tabular}{llrrrrrrrr}
    \toprule
    & & \multicolumn{4}{c}{CIFAR-10}
      & \multicolumn{4}{c}{CIFAR-100} \\
    \cmidrule(lr){3-6}\cmidrule(lr){7-10}
    Method & T. & Clean & AA & Avg & NRR
           & Clean & AA & Avg & NRR \\
    \midrule
    PGD-AT & -- & 84.54 & 43.57 & 64.06 & 57.50
                 & 57.36 & 20.34 & 38.85 & 30.03 \\
    TRADES & -- & 81.71 & 49.16 & 65.44 & 61.39
                 & 55.28 & 23.55 & 39.42 & 33.03 \\
    MART   & -- & 80.41 & 45.56 & 62.99 & 58.16
                 & 53.86 & 22.72 & 38.29 & 31.96 \\
    \midrule
    ARD & nat. & 84.58 & 43.75 & 64.16 & 57.67
               & 57.61 & 20.24 & 38.93 & 29.96 \\
    RSLAD & nat. & 85.15 & 42.91 & 64.03 & 57.06
                 & 59.68 & 21.30 & 40.49 & 31.40 \\
    AdaAD & nat. & 85.39 & 46.22 & 65.81 & 59.98
                 & 59.79 & 23.19 & 41.49 & 33.42 \\
    AdaAD + IGDM & nat. & 85.01 & 45.50 & 65.26 & 59.27
                        & 48.25 & 19.48 & 33.87 & 27.75 \\
    LBGAT & nat. & -- & -- & -- & --
                 & 56.46 & 26.02 & 41.24 & 35.62 \\
    ARREST & nat. & 85.06 & 46.42 & 65.74 & 60.06
                  & \textbf{66.97} & 22.78 & 44.88 & 34.00 \\
    \midrule
    HAT & nat. & 84.59 & 48.81 & 66.70 & 61.90
               & 60.67 & 24.45 & 42.56 & 34.85 \\
    ADR & self & 83.36 & 49.72 & 66.54 & 62.29
               & 59.93 & 27.20 & 43.57 & 37.42 \\
    RPAT++ & -- & 82.41 & 50.75 & 66.58 & 62.82
                  & 55.93 & 27.36 & 41.64 & 36.74 \\
    \midrule
    \textbf{CFA (ours)}
      & nat.
      & \textbf{87.37} & \textbf{51.33}
      & \textbf{69.35} & \textbf{64.67}
      & \textbf{64.82} & \textbf{28.04}
      & \textbf{46.43} & \textbf{39.15} \\
    \bottomrule
  \end{tabular}
\end{table}

\Cref{tab:main} compares CFA with standard adversarial training and robust
distillation methods on CIFAR-10 and CIFAR-100.
On CIFAR-100, CFA achieves \(64.82\%\) clean accuracy and \(28.04\%\) AA,
yielding an NRR of \(39.15\).
Compared with ADR, the closest baseline in NRR, CFA improves clean accuracy
by \(4.89\) points, AA by \(0.84\) points, and NRR by \(1.73\) points.
ARREST achieves higher clean accuracy at \(66.97\%\), but its AA is
\(5.26\) points lower than CFA, reflecting a different clean--robust
operating point.

On CIFAR-10, CFA reaches \(87.37\%\) clean accuracy and \(51.33\%\) AA,
with an NRR of \(64.67\).
Its NRR exceeds the closest baseline by \(1.85\) points.
Together, these results show that CFA maintains high clean accuracy while
remaining competitive in robust accuracy across both CIFAR datasets.


\subsection{Generalization Across Datasets and Architectures}

We next examine whether the behavior observed on the CIFAR benchmarks
extends to a larger-scale dataset and a higher-capacity architecture.

\begin{table}[!htbp]
  \centering
  \footnotesize
  \setlength{\tabcolsep}{3.2pt}
  \caption{
  Results on Tiny-ImageNet-200 with ResNet-18 and CIFAR-10 with
  WideResNet-34-10, measured in our framework.
  }
  \label{tab:scale}
  \begin{tabular}{llrrrrrrrr}
    \toprule
    & & \multicolumn{4}{c}{Tiny-ImageNet, ResNet-18}
      & \multicolumn{4}{c}{CIFAR-10, WRN-34-10} \\
    \cmidrule(lr){3-6}\cmidrule(lr){7-10}
    Method & T. & Clean & AA & Avg & NRR
           & Clean & AA & Avg & NRR \\
    \midrule
    PGD-AT & -- & 49.72 & 14.45 & 32.09 & 22.39
                 & 86.00 & 45.71 & 65.86 & 59.69 \\
    TRADES & -- & 46.98 & 15.95 & 31.46 & 23.83
                 & 85.13 & 48.99 & 67.06 & 62.19 \\
    MART & -- & 46.01 & 16.90 & 31.46 & 24.73
               & 85.29 & 47.06 & 66.18 & 60.65 \\
    \midrule
    ARD & nat. & 49.74 & 14.48 & 32.11 & 22.43
               & 86.35 & 45.52 & 65.94 & 59.61 \\
    RSLAD & nat. & 51.26 & 17.02 & 34.14 & 25.56
                 & 86.41 & 45.80 & 66.10 & 59.87 \\
    AdaAD & nat. & 51.59 & 16.79 & 34.19 & 25.33
                 & 86.80 & 46.84 & 66.82 & 60.85 \\
    LBGAT & nat. & 51.04 & 18.47 & 34.75 & 27.13
                 & 81.56 & 51.79 & 66.68 & 63.35 \\
    ARREST & nat. & \textbf{59.38} & 16.06 & 37.72 & 25.28
                  & 87.54 & 48.83 & 68.19 & 62.69 \\
    HAT & nat. & 52.57 & 17.28 & 34.93 & 26.01
               & 87.15 & 53.21 & 70.18 & 66.07 \\
    ADR & self & 47.66 & \textbf{20.85} & 34.25 & 29.01
               & 88.64 & 52.34 & 70.49 & 65.82 \\
    \midrule
    \textbf{CFA (ours)}
      & nat.
      & 55.16 & 20.54 & \textbf{37.85} & \textbf{29.93}
      & \textbf{88.67} & \textbf{55.29}
      & \textbf{71.98} & \textbf{68.11} \\
    \bottomrule
  \end{tabular}
\end{table}

On Tiny-ImageNet-200, CFA achieves the highest Avg and NRR without maximizing
either endpoint individually (\Cref{tab:scale}).
ARREST attains higher clean accuracy, while ADR exceeds CFA in AA by
\(0.31\) points.
CFA instead reaches \(55.16\%\) clean accuracy and \(20.54\%\) AA,
resulting in the highest NRR of \(29.93\).

The effect of the clean self-teacher observed in
\Cref{sec:teacher_transfer} is also present on Tiny-ImageNet.
Replacing the 200-epoch reference teacher with an 80-epoch teacher changes
the resulting student from \(55.16/20.54\) to \(57.08/18.96\)
clean/AA.
This again illustrates that changing the clean teacher can shift the
student's clean--robust operating point rather than uniformly improving both
metrics.

With WideResNet-34-10 on CIFAR-10, CFA achieves \(88.67\%\) clean accuracy
and \(55.29\%\) AA.
Its clean accuracy is nearly identical to ADR
(\(88.64\%\)), while its AA is \(2.95\) points higher.
On CIFAR-100 with the same architecture, CFA reaches
\(65.04/30.14\) clean/AA, compared with \(64.50/29.04\) for ADR
(\Cref{tab:wrn100}).
These results indicate that the observed behavior is not restricted to
ResNet-18 or to a single dataset.


\subsection{Ablation and Component Analysis}
\label{sec:components}

We next examine the components underlying CFA.
We first isolate the clean feature-anchoring objective, then compare feature
and logit targets, and finally study the contribution and assignment of
sample-dependent perturbation radii.
We separately report the effects of weight averaging and AWP.


\paragraph{Clean feature anchoring.}

We first test whether the gain arises from the clean feature anchor rather
than from the surrounding adversarial-training recipe.
Under the same base training regime, replacing feature anchoring with
label-based cross-entropy reduces both clean and robust accuracy.

\begin{table}[!htbp]
  \centering
  \small
  \setlength{\tabcolsep}{5pt}
  \caption{
  Objective control on CIFAR-100 with ResNet-18, 100 epochs, no weight
  averaging or AWP, \(p=0\), and
  \(\epsilon_{\mathrm{train}}=8.8/255\).
  }
  \label{tab:controls}
  \begin{tabular}{lrrr}
    \toprule
    Objective & Clean & AA & NRR \\
    \midrule
    Anchor
      & \textbf{61.29} & \textbf{25.36} & \textbf{35.88} \\
    Label cross-entropy, same recipe
      & 54.60 & 19.45 & 28.68 \\
    Label cross-entropy, own recipe
      & 56.66 & 20.91 & 30.55 \\
    \bottomrule
  \end{tabular}
\end{table}

As shown in \Cref{tab:controls}, replacing the anchor objective with
cross-entropy under the same recipe reduces clean accuracy by \(6.69\)
points and AA by \(5.91\) points.
Allowing cross-entropy to use its own optimizer and schedule narrows the gap,
but CFA remains ahead by \(4.63\) clean and \(4.45\) AA points.
Initialization and schedule controls change NRR by at most \(0.22\)
(\Cref{tab:recipe}), indicating that the difference cannot be attributed
solely to these training choices.


\paragraph{Feature and logit targets.}

The analysis in \Cref{sec:fragile} showed that fixed-target regression need
not be restricted to feature coordinates.
We therefore compare feature anchoring with direct logit regression under
otherwise comparable training conditions.

At a uniform radius of \(8.8/255\) and 100 training epochs, logit regression
reaches \(60.36/28.02\) clean/AA, compared with \(60.42/28.42\) for
feature regression (\Cref{tab:regression_stack}).
Under the final 150-epoch recipe with sensitivity-matched radii, logit
regression reaches \(64.49/27.85\) clean/AA (NRR \(38.90\)), compared with
\(64.82/28.04\) (NRR \(39.15\)) for feature regression
(\Cref{tab:finalallocation}).
The small differences are consistent with our analysis: effective transfer
through a fixed clean target is not confined to feature space.


\paragraph{Sensitivity-matched radius allocation.}

We next isolate the second component of CFA: allocating the perturbation
radius according to local sensitivity.
We hold the teacher, training schedule, mean radius, weight averaging, and
AWP fixed and vary only the radius allocation.

\begin{table}[!htbp]
  \centering
  \small
  \setlength{\tabcolsep}{5pt}
  \caption{
  Clean-target regression in the final CIFAR-100 recipe
  (ResNet-18, 150 epochs, \(\epsilon_{\mathrm{train}}=8/255\)).
  Feature rows report mean \(\pm\) sample standard deviation over three
  student seeds and differ only in \(p\); logit regression is one run at
  \(p=1\).
  }
  \label{tab:finalallocation}
  \begin{tabular}{lrrr}
    \toprule
    Target and radius & Clean & AA & NRR \\
    \midrule
    Features, uniform (\(p=0\))
      & \(62.63\pm0.22\)
      & \(28.01\pm0.20\)
      & \(38.71\pm0.21\) \\
    Features, sensitivity-matched (\(p=1\))
      & \(\mathbf{64.82\pm0.07}\)
      & \(\mathbf{28.04\pm0.09}\)
      & \(\mathbf{39.15\pm0.09}\) \\
    Logits, sensitivity-matched (\(p=1\))
      & 64.49 & 27.85 & 38.90 \\
    \bottomrule
  \end{tabular}
\end{table}

Replacing the uniform radius with sensitivity matching increases clean
accuracy by \(2.19\pm0.20\) points across three paired runs, while AA changes
by only \(0.03\pm0.23\) points (\Cref{tab:finalallocation}).
The clean-accuracy improvement appears in every seed, whereas the AA
difference is comparable to run-to-run variation.
Thus, in the final training configuration, sensitivity matching primarily
shifts the operating point toward higher clean accuracy without an observed
loss in AA.
This empirical behavior is not implied by the first-order motivation in
\Cref{sec:angeps}, which only motivates how the perturbation budget is
allocated.


\paragraph{Which signal should allocate the radius?}

The previous comparison establishes the effect of sensitivity matching
relative to a uniform radius, but does not determine whether the sensitivity
signal itself matters or whether heterogeneous radii are sufficient.
We therefore compare sensitivity with alternative assignments under the same
mean perturbation budget.

\begin{table}[!htbp]
  \centering
  \small
  \setlength{\tabcolsep}{4.5pt}
  \caption{
  Radius-allocation signals under the same mean radius on CIFAR-100 with
  ResNet-18, trained for 100 epochs with WA and AWP.
  \(\mathrm{sd}(w)\) measures multiplier dispersion.
  }
  \label{tab:allocation}
  \begin{tabular}{lrrrr}
    \toprule
    Allocation & \(\mathrm{sd}(w)\) & Clean & AA & NRR \\
    \midrule
    Uniform
      & 0 & 60.42 & 28.42 & 38.66 \\
    Sensitivity
      & 0.32 & \textbf{62.17} & \textbf{28.86} & \textbf{39.42} \\
    Same multipliers, difficulty order
      & 0.32 & 60.90 & 28.06 & 38.42 \\
    Difficulty magnitude
      & 0.39 & 61.27 & 28.12 & 38.55 \\
    Logit margin
      & 0.42 & 61.07 & 27.79 & 38.20 \\
    \bottomrule
  \end{tabular}
\end{table}

Sensitivity matching gives the highest clean accuracy, AA, and NRR among the
tested allocation signals (\Cref{tab:allocation}).
In particular, reassigning the same set of sensitivity-derived radius
multipliers according to sample difficulty reduces AA by \(0.80\) points and
NRR from \(39.42\) to \(38.42\).
Because this control preserves the radius distribution while changing which
example receives each radius, the result indicates that the gain cannot be
explained solely by introducing heterogeneous perturbation radii.
Difficulty magnitude and logit-margin allocation also yield lower NRR than
the sensitivity-based assignment.


\paragraph{Weight averaging and AWP.}

Finally, we examine the additional training components used in the final
configuration.
Weight averaging and AWP are not defining components of the CFA objective,
but are included in the final recipe to strengthen adversarial training.
In the ordered ablation, adding weight averaging and AWP increases AA by
\(1.95\) and \(1.17\) points, respectively
(\Cref{tab:ladder}).
Because this ordered ablation does not isolate interactions between the two
components, additional radius and training-stack comparisons are reported in
\Cref{tab:radius,tab:stackboth}.

To verify that the final training stack alone does not account for CFA's
performance, we also apply teacher initialization, weight averaging, and AWP
to the baseline methods.

\begin{table}[!htbp]
  \centering
  \scriptsize
  \setlength{\tabcolsep}{4pt}
  \caption{
  Adding teacher initialization, weight averaging, and AWP to baseline
  recipes on CIFAR-100 (ResNet-18).
  }
  \label{tab:baselinestack}
  \begin{tabular}{lrrrcrrr}
    \toprule
    & \multicolumn{3}{c}{Own recipe}
    && \multicolumn{3}{c}{+ teacher init + WA + AWP} \\
    \cmidrule(lr){2-4}\cmidrule(lr){6-8}
    Method & Clean & AA & NRR && Clean & AA & NRR \\
    \midrule
    PGD-AT & 57.36 & 20.34 & 30.03
           && 59.85 & 25.77 & 36.03 \\
    TRADES & 55.28 & 23.55 & 33.03
           && 56.16 & 24.56 & 34.17 \\
    MART & 53.86 & 22.72 & 31.96
         && 50.97 & 25.13 & 33.66 \\
    ARD & 57.61 & 20.24 & 29.96
        && 60.19 & 26.10 & 36.41 \\
    RSLAD & 59.68 & 21.30 & 31.40
          && 63.44 & 25.55 & 36.43 \\
    AdaAD & 59.79 & 23.19 & 33.42
          && 57.21 & 27.10 & 36.78 \\
    AdaAD + IGDM & 48.25 & 19.48 & 27.75
                 && 54.24 & 25.56 & 34.75 \\
    Consistency-AT & 57.53 & 20.01 & 29.69
                   && 60.74 & 25.38 & 35.80 \\
    LBGAT & 56.46 & 26.02 & 35.62
          && 53.91 & 25.16 & 34.31 \\
    ADR & 59.93 & 27.20 & 37.42
        && 59.84 & 27.84 & 38.00 \\
    HAT & 60.67 & 24.45 & 34.85
        && 57.63 & 24.74 & 34.62 \\
    ARREST & 66.44 & 22.99 & 34.16
           && 64.95 & 23.02 & 33.99 \\
    \midrule
    \textbf{CFA (ours)}
      & \multicolumn{3}{c}{---}
      && \textbf{64.82} & \textbf{28.04} & \textbf{39.15} \\
    \bottomrule
  \end{tabular}
\end{table}

The additional training components increase NRR for nine of the twelve
baselines, but none reaches CFA's NRR of \(39.15\)
(\Cref{tab:baselinestack}); ADR is the closest at \(38.00\).
This comparison indicates that the performance of CFA cannot be attributed
solely to teacher initialization, weight averaging, or AWP.
ARREST uses its 20-epoch recipe in this control, whereas
\Cref{tab:main} reports its 150-epoch comparison; ADR already includes a
SAM-style weight perturbation in its original recipe.

Attack diagnostics, repeat-seed results, and additional implementation
controls are provided in \Cref{app:gradmask,app:repro}.